\documentclass[11pt,a4paper]{amsart}

\usepackage[utf8]{inputenc}
\usepackage[T1]{fontenc}
\usepackage[ngerman]{babel}
\usepackage{amsmath,amssymb,amsthm}
\usepackage{array,booktabs,longtable}
\usepackage{hyperref}

\newtheorem{theorem}{Satz}[section]
\newtheorem{proposition}[theorem]{Proposition}
\newtheorem{lemma}[theorem]{Lemma}
\newtheorem{corollary}[theorem]{Korollar}
\newtheorem{definition}[theorem]{Definition}
\newtheorem{remark}[theorem]{Bemerkung}
\newtheorem{openproblem}[theorem]{Offenes Problem}

\newcommand{\RH}{\mathrm{RH}}
\newcommand{\Tr}{\operatorname{Tr}}
\newcommand{\FP}{\operatorname{FP}}
\newcommand{\AC}{\operatorname{AC}}

\title[P08 --- Renormalized Prime Operators]{P08 --- Renormalized Prime Operators and Finite-Part Structures}
\author{Objekt-X-Programm}
\date{2026-08-09}

\begin{document}

\begin{abstract}
P08 konsolidiert zwei logisch getrennte Stränge des Grenzoperator- und
Renormierungsprogramms. Der erste ist eine Renormierungsdiagnose für endliche
Jacobi-Modelle: die erste Lanczos-Kante kollabiert streng und eine mögliche
Divergenz von $b_{2,N}/b_{1,N}$ würde eine skalare Renormierung ausschließen,
ohne eine nichtskalare Prä-Lanczos-Metrik zu konstruieren. Der zweite Strang
untersucht renormalisierte Primoperatoren, Mangoldt-Spuren und
Finite-Part-Strukturen. Eine feste-$\beta$-Spurklasse ist nur konditional im
relativen Rang-eins-Modell verfügbar. Der exakte Mellin-Kanal ist die
glättete Mangoldt-Summe $\Psi_{\varphi,X}$, nicht die Prime-only-Summe mit
Cutoff $\varphi(p/X)$. Die operatorielle Identifikation eines konkreten
Finite Parts mit der analytischen Fortsetzung von $-\zeta'/\zeta$ bleibt
offen. Es wird weder Objekt X noch ein Hilbert--Pólya-Endoperator konstruiert,
und kein RH-Beweis behauptet.
\end{abstract}

\maketitle

% ----------------------------------------------------------------
\section{Bindende Typ- und Statusfirewalls}
% ----------------------------------------------------------------

\begin{remark}[P05-Scope]\label{rem:p05-scope}
Für P08 gilt
\[
  \operatorname{rank}C_p^{\rm rel}\le1
\]
nur in der induzierten relativen Modellrealisierung. Der gewichtete
Rang-eins-Operator ist
\[
  P_p=|c_p|^2\Pi_p^{(1)},
\]
und im Allgemeinen kein orthogonaler Projektor. Nichtentartung $c_p\ne0$,
Liftunabhängigkeit von $|c_p|^2$ und eine intrinsische termweise Asymptotik
bleiben offen.
\end{remark}

\begin{remark}[Zwei Jacobi-Typen]\label{rem:jacobi-types}
Die endliche direkt symmetrische Schließung
\[
  A_N^{\rm sym}=B_N^\Lambda=J_N^\Lambda+(J_N^\Lambda)^*
\]
ist selbstadjungiert. Davon ist der historische antisymmetrische Pfad zu
trennen:
\[
  J_N^-:=\frac12(\Theta_N-\Theta_N^\dagger),
  \qquad
  S_N:=\frac1{2i}(\Theta_N-\Theta_N^\dagger)=-iJ_N^-.
\]
$J_N^-$ ist schiefadjungiert, $S_N$ selbstadjungiert. Eine konkrete
selbstadjungierte Realisierung des historischen $A_N^{\rm Jac,-}$ bleibt
offen.
\end{remark}

\begin{definition}[Analytische Regularisierungsdefinition]
\label{def:trreg}
Die Schreibweise
\[
  \Tr_{\rm reg}(R\Sigma)(\beta)
  :=\AC\!\left[-\frac{\zeta'}{\zeta}\right](\beta)
\]
ist eine Definition. Sie ist keine operatorielle Herleitung eines
Cutoff-Grenzwerts.
\end{definition}

% ----------------------------------------------------------------
\section{Moment-/KMS-Eingang und Herglotz-Grenze}
% ----------------------------------------------------------------

\begin{theorem}[Korrigierte $\xi$-Konstante]\label{thm:cxi}
Für die in NEU-121 verwendete Größe gilt
\[
  C_\xi=-\frac{\xi'(0)}{\xi(0)}
  =1+\frac{\gamma_E}{2}-\frac12\log(4\pi)
  \approx0.0230957.
\]
Der historische Zahlenwert $-0.5493$ ist gesperrt und superseded.
\end{theorem}

\begin{theorem}[KMS/GNS-Normalisierung]\label{thm:gns-normalization}
Sei
\[
  \tau_{\beta,N}(T)=\sum_{n\le N}n^{-\beta}\langle e_n,Te_n\rangle,
  \qquad Z_{\beta,N}=\tau_{\beta,N}(1),
\]
und $\varphi_{\beta,N}=Z_{\beta,N}^{-1}\tau_{\beta,N}$. Dann lautet die
korrigierte GNS-Normalisierung
\[
  \widehat\Omega=Z^{-1/2}\Omega_\tau.
\]
Für $\beta=1$ gilt
\[
  Z_{1,N}^{-1}\sim\frac1{\log N}.
\]
Die historischen Formeln $Z^{+1/2}\Omega_\tau$ und $R_N\sim\log N$ werden
nicht übernommen.
\end{theorem}

\begin{theorem}[Herglotz-Firewall]\label{thm:herglotz-firewall}
P07 fixiert
\[
  m_{\rm arith}\text{ ist Herglotz}\quad\Longleftrightarrow\quad\RH.
\]
Daher gilt für echte Herglotz-Approximanten die logische Implikation
\[
  \widetilde m_N^{\rm ren}
  \xrightarrow{\rm loc.\,unif.}m_{\rm arith}
  \Longrightarrow
  m_{\rm arith}\text{ Herglotz}
  \Longrightarrow\RH.
\]
Die Existenz und kanonische Wahl einer solchen Approximantenfolge wird damit
nicht bewiesen.
\end{theorem}

\begin{openproblem}[KMS/GNS--Jacobi-Brücke]\label{op:kms-jacobi}
Offen bleiben eine vorab fixierte selbstadjungierte GNS-Observable mit der
gewünschten Lanczos-Darstellung, die Formkompatibilität mit der
Bombieri-/Weil-Geometrie, die konkrete Selbstadjungiertheit des historischen
$A_N^{\rm Jac,-}$ sowie die kanonische Nevanlinna-Normalisierung und
Tail-Kontrolle.
\end{openproblem}

% ----------------------------------------------------------------
\section{Renormierungsdiagnose des symmetrischen Jacobi-Pfads}
% ----------------------------------------------------------------

\begin{theorem}[Kollaps der ersten Lanczos-Kante]\label{thm:b1-collapse}
Für den direkt symmetrisierten endlichen Kandidaten gilt
\[
  a_{0,N}=0,
\]
und
\[
  b_{1,N}
  =\frac{\gamma}{N}\sqrt{\sum_{n=2}^{N-1}\Lambda(n)^2}
  \asymp\gamma\sqrt{\frac{\log N}{N}}
  \longrightarrow0.
\]
Damit scheitert die unrenormierte nichtdegenerierte Eintrittsbedingung
$b_{1,N}\to b_1>0$.
\end{theorem}

\begin{corollary}[Startvektor-Weylkanal]\label{cor:start-weyl}
Aus $A_N^{\rm sym}e_0=b_{1,N}e_1$ folgt für
$z\in\mathbb C\setminus\mathbb R$
\[
  \langle e_0,(A_N^{\rm sym}-z)^{-1}e_0\rangle
  \longrightarrow-\frac1z.
\]
Der unrenormierte Startvektor-Kanal kann daher nicht $m_{\rm arith}$ liefern.
\end{corollary}

\begin{remark}[Keine globale Diagonalität]\label{rem:no-global-diag}
Aus $b_{1,N}\to0$ folgt nicht, dass alle höheren Offdiagonalparameter
verschwinden oder ein Grenzoperator insgesamt diagonal wird. Diese historische
Folgerung ist ein No-Go.
\end{remark}

\begin{proposition}[Skalare Lanczos-Kovarianz]\label{prop:lanczos-cov}
Für $c_N>0$ gilt
\[
  A_N\mapsto c_NA_N
  \Longrightarrow
  a_{j,N}\mapsto c_Na_{j,N},\qquad
  b_{j,N}\mapsto c_Nb_{j,N}.
\]
Insbesondere ist $b_{2,N}/b_{1,N}$ unter positiver skalarer
Prä-Lanczos-Skalierung invariant.
\end{proposition}

\begin{proposition}[Abstraktes No-scalar-Lemma]\label{prop:no-scalar}
Falls
\[
  \frac{b_{2,N}}{b_{1,N}}\longrightarrow\infty,
\]
kann keine positive skalare Folge $\kappa_N$ beide Größen
$b_{1,N}/\kappa_N$ und $b_{2,N}/\kappa_N$ gegen endliche positive Grenzwerte
schicken.
\end{proposition}

\begin{openproblem}[Zweite Kante und Prä-Lanczos-Geometrie]
\label{op:wn}
Die endlichen Daten liefern starke Evidenz für Wachstum von
$b_{2,N}/b_{1,N}$, beweisen aber den Grenzwert gegen $\infty$ nicht. Ebenso
bleibt die Konstruktion einer intrinsischen positiven nichtskalaren
Prä-Lanczos-Geometrie $W_N$ offen. Die Barriere diagnostiziert einen möglichen
Bedarf; sie konstruiert $W_N$ nicht.
\end{openproblem}

% ----------------------------------------------------------------
\section{Prä-Lanczos-/Grammetrik und lokale No-Gos}
% ----------------------------------------------------------------

\begin{remark}[Ebenentrennung]\label{rem:levels}
Eine Jacobi-seitige Self-Energy ist noch keine intrinsische
Prä-Lanczos-Metrik. Ein Kandidat $W_N$ müsste vor Lanczos auf der
Quell-/Feshbach-Ebene konstruiert werden.
\end{remark}

\begin{proposition}[Korrekte Rang-eins-Form im Modell]\label{prop:rank-one-form}
Im historischen Rang-eins-Modell lautet die Operatorform
\[
  \Sigma_N(\beta)x
  =\sum_p w_p(\beta)\Psi_p\langle\Psi_p,x\rangle,
\]
beziehungsweise die Quadratform
\[
  \langle x,\Sigma_N(\beta)x\rangle
  =\sum_p w_p(\beta)|\langle\Psi_p,x\rangle|^2.
\]
Die historische Vektor-gleich-Skalar-Formel wird nicht übernommen.
\end{proposition}

\begin{remark}[Paper-VII-Skalierungs-No-Go]\label{rem:h3-nogo}
Aus $P_{kl}\le C_2c^{1/2}$ und $A_{kl}=P_{kl}c^{1/2}$ folgt
$A_{kl}=O(c)$, nicht $O(1)$. Die historische H3-Verifikation mit dieser
Skalierung ist gesperrt.
\end{remark}

\begin{remark}[Cancellation-No-Go]\label{rem:cancellation-nogo}
Phasencancellation signierter dyadischer Summen kontrolliert keine absolute
Schur-Zeilensumme. Das Modell
\[
  T_{ij}=\frac{e^{i\alpha(i-j)}}{|i-j|}
\]
weist Cancellation auf, während
\[
  \sum_{j\ne i}|T_{ij}|\asymp\log N.
\]
Eine korrekte Nelson-/Schur-Brücke benötigt daher einen separaten
Operatornorm-, Quadratsummen-, $TT^*$-, Cotlar- oder Orthogonalitätsmechanismus.
\end{remark}

\begin{remark}[PSWF-Route]\label{rem:pswf}
Die PSWF-/Edge-Koerzivitätsidee bleibt Methodenheuristik. Insbesondere wird
keine Operatoridentifikation $D_{\rm rel}=\overline{iJ^-}$ und keine konkrete
Prä-Lanczos-Metrik daraus abgeleitet.
\end{remark}

% ----------------------------------------------------------------
\section{Renormalisierte Selbstenergie und feste-$\beta$-Spurklasse}
% ----------------------------------------------------------------

\begin{definition}[Algebraische $\beta$-Zerlegung]\label{def:beta-split}
Aus
\[
  \frac1{1-p^{-\beta}}
  =1+\frac{p^{-\beta}}{1-p^{-\beta}}
\]
folgt formal
\[
  \Sigma_{\rm rel}(\beta)
  =\Sigma_{\rm rel}^{\infty}+\Sigma_{\rm rel}^{\rm ren}(\beta),
\]
mit
\[
  \Sigma_{\rm rel}^{\rm ren}(\beta)
  :=\sum_p\frac{p^{-\beta}}{1-p^{-\beta}}P_p.
\]
\end{definition}

\begin{remark}[Rohanteil-Firewall]\label{rem:raw-firewall}
Aus den vorhandenen Upper Bounds folgt weder die Divergenz noch die
Nicht-Regularisierbarkeit von $\Sigma_{\rm rel}^{\infty}$. Die korrekte
PNT-Vergleichsasymptotik lautet
\[
  \sum_{p\le N}\frac{(\log p)^2}{p}
  \sim\frac12(\log N)^2,
\]
nicht $\frac13(\log N)^3$.
\end{remark}

\begin{proposition}[Konditionale feste-$\beta$-Spurklasse]
\label{prop:s1-fixed-beta}
Unter
\[
  \operatorname{rank}C_p^{\rm rel}\le1
  \quad\text{im relativen Modell}
\]
und
\[
  |c_p|^2=O\!\left(\frac{(\log p)^2}{p}\right)
\]
folgt für jedes feste reelle $\beta>0$
\[
  \Sigma_{\rm rel}^{\rm ren}(\beta)\in\mathcal S_1,
\]
gleichmäßig für $\beta\ge\beta_0>0$. Dies ist ein konditionaler
modellrelativer Satz. T2 und $c_p\ne0$ sind keine Voraussetzungen dieses
reinen $\mathcal S_1$-Schritts.
\end{proposition}

\begin{openproblem}[Quantitative Kanalnorm]\label{op:cp-bound}
Die Schranke
\[
  |c_p|^2=O\!\left(\frac{(\log p)^2}{p}\right)
\]
ist intrinsisch nicht bewiesen. Ihr Modellursprung enthält insbesondere den
offenen Schritt $B_p=O(1/p)$.
\end{openproblem}

\begin{proposition}[Fredholm-Basistheorie]\label{prop:fredholm-base}
Sobald $\Sigma_{\rm rel}^{\rm ren}(\beta)\in\mathcal S_1$ vorliegt, sind die
gewöhnliche Fredholm-Determinante und die Potenzspuren im
Standard-Schattenrahmen definiert. Ohne T2 sind die primweisen Gewichte jedoch
nicht automatisch Eigenwerte der Gesamtsumme; ein reines Euler-/Ihara-Produkt
folgt nicht aus der Spurklasse allein.
\end{proposition}

\begin{proposition}[Korrekte zweite Spur]\label{prop:trace2}
Im endlichen beziehungsweise spurklassigen Scope gilt formal
\[
  \Tr(\Sigma^2)=\sum_{p,q}w_pw_q\Tr(P_pP_q).
\]
Die historische Formel mit einem zusätzlichen Gesamtfaktor wird nicht
übernommen.
\end{proposition}

% ----------------------------------------------------------------
\section{T2, Nichtentartung und Mangoldt-Observable}
% ----------------------------------------------------------------

\begin{lemma}[Konditionale Edge-Label-Orthogonalität]\label{lem:t2}
Falls eine intrinsisch gerechtfertigte orthogonale Zerlegung
\[
  W_{\rm res,rel}=\bigoplus_{(m,p)}^\perp H_{m\to pm}
\]
vorliegt, sind verschiedene Primlabelkanäle orthogonal. Die Herleitung dieser
orthogonalen Summe aus der ursprünglichen BC-Geometrie ist nicht geschlossen;
intrinsisches T2 bleibt offen.
\end{lemma}

\begin{openproblem}[Nichtentartung]\label{op:nonzero-cp}
Ein allgemeiner Beweis
\[
  c_p\ne0\qquad\text{für alle relevanten Primzahlen}
\]
liegt nicht vor.
\end{openproblem}

\begin{definition}[Primdiagonale Mangoldt-Observable, conditional]
\label{def:mangoldt-r}
Unter T2, $c_p\ne0$ und einer abgeschlossenen primdiagonalen
Hilbertraumrealisierung kann formal
\[
  R_p^{\rm Mang}:=\frac{\log p}{|c_p|^2}
\]
definiert werden. Im modellrelativen Rang-eins-Scope gilt dann
\[
  \Tr(RP_p)=\log p.
\]
Eine konkrete selbstadjungierte Operatorrealisierung von $R$ bleibt offen.
Unter der zusätzlichen konditionalen Kanalnormschranke folgt lediglich
\[
  R_p^{\rm Mang}\gtrsim\frac{p}{\log p};
\]
eine Asymptotik ist nicht bewiesen.
\end{definition}

\begin{proposition}[Gewöhnliche Mangoldt-Spur]\label{prop:mangoldt-trace}
Im primdiagonalen Modell gilt für $\Re\beta>1$ formal
\[
  \Tr\bigl(R\Sigma_{\rm rel}^{\rm ren}(\beta)\bigr)
  =\sum_p\frac{\log p\,p^{-\beta}}{1-p^{-\beta}}
  =-\frac{\zeta'}{\zeta}(\beta).
\]
Dies ist ein konditionaler Modellsatz im gewöhnlichen absoluten
Konvergenzgebiet und keine Finite-Part-Aussage.
\end{proposition}

% ----------------------------------------------------------------
\section{Scharfer Primcutoff und Prime-only-Firewall}
% ----------------------------------------------------------------

\begin{definition}[Prime-only-Cutoff]\label{def:sharp-prime-cutoff}
\[
  S_X(\beta)
  :=\sum_{p\le X}\frac{\log p\,p^{-\beta}}{1-p^{-\beta}}
  =\sum_{k\ge1}T_k(X,\beta),
\]
\[
  T_k(X,\beta):=\sum_{p\le X}\log p\,p^{-k\beta}.
\]
\end{definition}

\begin{theorem}[PNT-Hauptterm]\label{thm:prime-layer-main}
Für festes $k\beta$ mit $\Re(k\beta)<1$ gilt
\[
  T_k(X,\beta)
  \sim\frac{X^{1-k\beta}}{1-k\beta}.
\]
Für $k\beta=1$ tritt $\log X$ auf. Am Rand $\Re(k\beta)=1$,
$k\beta\ne1$, ist der Hauptterm beschränkt oszillierend und besitzt im
Allgemeinen keinen Grenzwert.
\end{theorem}

\begin{remark}[$\vartheta$ ist nicht $\psi$]\label{rem:theta-psi}
$T_k$ ist eine Prime-only-Summe über
\[
  \vartheta(X)=\sum_{p\le X}\log p,
\]
während die klassische explizite Formel direkt
\[
  \psi(X)=\sum_{p^j\le X}\log p
\]
beschreibt. Ein direkter Import der $\psi$-Nullstellenterme in $T_k$ ohne
Primpotenz-/Möbiuskorrektur ist gesperrt. Die RH-Verbindung bleibt nur als
partielle Richtungsstruktur erhalten; ein vollständiger Prime-only-
Äquivalenzsatz ist offen.
\end{remark}

% ----------------------------------------------------------------
\section{Der korrekte Mellin-Kanal}
% ----------------------------------------------------------------

\begin{definition}[Testfunktion und Mellintransformierte]
\label{def:mellin-test}
Sei $\varphi\in C_c^\infty([0,\infty))$ mit $\varphi(x)=1$ nahe $0$. Für
$\Re s>0$ sei
\[
  \widehat\varphi(s)
  :=\int_0^\infty\varphi(x)x^{s-1}\,dx.
\]
Die Mellintransformierte besitzt eine meromorphe Fortsetzung mit einfachem
Pol bei $s=0$ und
\[
  \operatorname{Res}_{s=0}\widehat\varphi(s)=1.
\]
Sie wird weder als ganz noch durch einen endlichen Wert
$\widehat\varphi(0)=1$ normiert.
\end{definition}

\begin{definition}[Geglättete Mangoldt-Summe]\label{def:psi-smooth}
\[
  \Psi_{\varphi,X}(\beta)
  :=\sum_{n\ge1}\Lambda(n)\varphi(n/X)n^{-\beta}
  =\sum_{p,k\ge1}\log p\,\varphi(p^k/X)p^{-k\beta}.
\]
\end{definition}

\begin{theorem}[Exakte Mellin-Darstellung]\label{thm:psi-mellin}
Für
\[
  c>\max(0,1-\Re\beta)
\]
gilt
\[
  \Psi_{\varphi,X}(\beta)
  =\frac1{2\pi i}\int_{(c)}\widehat\varphi(s)X^s
  \left(-\frac{\zeta'}{\zeta}(\beta+s)\right)\,ds.
\]
\end{theorem}

\begin{remark}[Prime-only-Mellin-No-Go]\label{rem:s-mellin-nogo}
Für
\[
  S_{\varphi,X}(\beta)
  =\sum_p\varphi(p/X)\frac{\log p\,p^{-\beta}}{1-p^{-\beta}}
\]
entsteht nach Mellin-Inversion
\[
  \sum_{p,k\ge1}\log p\,p^{-k\beta-s},
\]
nicht
\[
  -\frac{\zeta'}{\zeta}(\beta+s)
  =\sum_{p,k\ge1}\log p\,p^{-k\beta-ks}.
\]
Die historische NEU-148.A-Identität ist deshalb gesperrt.
\end{remark}

\begin{proposition}[Korrekte $\Psi/S$-Differenz]\label{prop:psi-s-diff}
Algebraisch gilt
\[
  \Psi_{\varphi,X}(\beta)-S_{\varphi,X}(\beta)
  =\sum_{k\ge2}\sum_p\log p\,
  [\varphi(p^k/X)-\varphi(p/X)]p^{-k\beta}.
\]
Für $\Re\beta>1/2$ ist die höhere-Primpotenzreihe absolut summierbar. Der für
die Finite-Part-Anwendung benötigte quantitative/uniforme Transfer bleibt im
versiegelten P08-Endstand offen; hier wird keine neue Hochstufung vorgenommen.
\end{proposition}

% ----------------------------------------------------------------
\section{Konturrest und analytischer Finite Part}
% ----------------------------------------------------------------

\begin{proposition}[Mellin-Pol bei $s=0$]\label{prop:mellin-zero-pole}
Ist $\beta$ keine Polstelle von $-\zeta'/\zeta$, dann liefert der Pol von
$\widehat\varphi$ bei $s=0$ im korrigierten $\Psi$-Strang
\[
  \operatorname{Res}_{s=0}
  \left[
    \widehat\varphi(s)X^s
    \left(-\frac{\zeta'}{\zeta}(\beta+s)\right)
  \right]
  =-\frac{\zeta'}{\zeta}(\beta).
\]
\end{proposition}

\begin{proposition}[Fixed-contour Restlemma, conditional]
\label{prop:fixed-contour}
Auf einer fest gewählten linken Kontur $\Gamma_{-M}$ mit
$\Re s\le-M<0$ und quantitativ positivem Abstand von allen relevanten Polen
gilt $|X^s|\le X^{-M}$. Zusammen mit geeignetem vertikalem Mellin-Abfall und
einer passenden Wachstumskontrolle von $\zeta'/\zeta$ ergibt sich conditional
\[
  R_{\varphi,X}(\beta)=O(X^{-M}).
\]
\end{proposition}

\begin{openproblem}[Uniforme Kontur und Residuenzählung]
\label{op:uniform-contour}
Eine einheitliche nullstellenvermeidende Kontur für einen variierenden
kompakten $\beta$-Bereich sowie die vollständige uniforme Residuenzählung sind
nicht bewiesen. Für kontinuierliches $K$ ist
\[
  \{\omega-\beta:\beta\in K\}
\]
insbesondere nicht diskret.
\end{openproblem}

\begin{openproblem}[Analytischer Finite Part]\label{op:analytic-fp}
Nach Abzug aller nichtverschwindenden Residuenbeiträge ist das Ziel
\[
  \FP_{X\to\infty}^{\varphi}\Psi_{\varphi,X}(\beta)
  \stackrel?=-\frac{\zeta'}{\zeta}(\beta).
\]
Das fixed-contour Lemma ist ein konditionaler Baustein; der vollständige
uniforme Abschluss bleibt offen.
\end{openproblem}

% ----------------------------------------------------------------
\section{Operatorielle Primlabel-Brücke und Cutoff-Hierarchie}
% ----------------------------------------------------------------

\begin{definition}[Primlabel-Observable, conditional]\label{def:np}
Unter einer intrinsisch gerechtfertigten orthogonalen Primzerlegung und
Nichtentartung kann formal
\[
  N_{\mathbb P}\Psi_p=p\Psi_p
\]
definiert werden. Für eine orthogonale, nicht notwendig normierte Familie
lautet die natürliche Multiplikationsdomäne
\[
  \sum_p p^2|\xi_p|^2\|\Psi_p\|^2<\infty,
\]
wobei auch das orthogonale Komplement festzulegen ist.
\end{definition}

\begin{proposition}[Formale Primlabel-Spuridentität]\label{prop:np-trace}
Im primdiagonalen Modell gilt algebraisch
\[
  \Tr\!\left(
    \varphi(N_{\mathbb P}/X)R\Sigma_{\rm rel}^{\rm ren}(\beta)
  \right)
  =S_{\varphi,X}(\beta).
\]
Dies ist ein konditionaler modellrelativer Satz.
\end{proposition}

\begin{openproblem}[Primlabel-Finite-Part]\label{op:np-fp}
Offen ist
\[
  \FP_{X\to\infty}^{N_{\mathbb P}}
  \Tr\!\left(
    \varphi(N_{\mathbb P}/X)R\Sigma_{\rm rel}^{\rm ren}(\beta)
  \right)
  \stackrel?=-\frac{\zeta'}{\zeta}(\beta).
\]
Dazu fehlen insbesondere der $\Psi/S$-Transfer in der benötigten Stärke und
die uniforme analytische Rest-/Residuenkontrolle.
\end{openproblem}

\begin{openproblem}[$R$-Cutoff]\label{op:r-cutoff}
Ein Cutoff $\varphi(R/\Lambda)$ ist nicht automatisch äquivalent zum
Primlabel-Cutoff. Selbst die Zusatzannahme
\[
  R_p\asymp\frac{p}{\log p}
\]
würde nur grobe Skalenvergleichbarkeit liefern. Zur Identifikation von
Finite Parts ist eine stärkere quantitative Asymptotik samt Fehlerkontrolle
nötig.
\end{openproblem}

% ----------------------------------------------------------------
\section{Zwei getrennte P08-Stränge}
% ----------------------------------------------------------------

\begin{remark}[Strang A: Renormierungsdiagnose]\label{rem:strand-a}
Gesichert ist $b_{1,N}\to0$. Der strenge Grenzwert
$b_{2,N}/b_{1,N}\to\infty$ bleibt offen. Falls er bewiesen wird, greift das
abstrakte No-scalar-Lemma. Eine positive nichtskalare Prä-Lanczos-Geometrie
$W_N$ bleibt dennoch eine separate offene Konstruktion.
\end{remark}

\begin{remark}[Strang B: Mangoldt/Mellin-Operatorbrücke]
\label{rem:strand-b}
Die korrekte Reihenfolge ist:
\begin{enumerate}
  \item modellrelative Rangstruktur plus quantitative $c_p$-Obergrenze
        $\Rightarrow$ feste-$\beta$-$\mathcal S_1$ conditional;
  \item T2 plus Nichtentartung plus abgeschlossene Primzerlegung
        $\Rightarrow$ primdiagonales $R$ conditional;
  \item exakte $\Psi_{\varphi,X}$-Mellin-Darstellung;
  \item uniforme Rest-/Residuenzählung;
  \item $\Psi/S$-Transfer;
  \item operatorielle Primlabel-Brücke;
  \item separater Vergleich zum $R$-Cutoff.
\end{enumerate}
Keine Stufe darf durch Definition~\ref{def:trreg} ersetzt werden.
\end{remark}

% ----------------------------------------------------------------
\section{Statusmatrix}
% ----------------------------------------------------------------

\small
\begin{longtable}{>{\raggedright\arraybackslash}p{0.57\textwidth}
                  >{\raggedright\arraybackslash}p{0.25\textwidth}
                  >{\raggedright\arraybackslash}p{0.10\textwidth}}
\toprule
Aussage & Status & Quelle\\
\midrule
\endfirsthead
\toprule
Aussage & Status & Quelle\\
\midrule
\endhead
korrigiertes $C_\xi\approx0.0230957$ & PROVED & 121Cfix\\
alter $C_\xi$-Wert $-0.5493$ & NO-GO/SUPERSEDED & 121\\
$\widehat\Omega=Z^{-1/2}\Omega_\tau$ & PROVED & 122\\
$Z_{1,N}^{-1}\sim1/\log N$ & PROVED & 121/122\\
historisches $A_N^{\rm Jac,-}$ selbstadjungiert & OPEN & H-T1\\
KMS/GNS $\to$ Jacobi / Formkompatibilität & OPEN & 122\\
$A_N^{\rm sym}$ endlich selbstadjungiert & PROVED & 123A\\
$b_{1,N}\asymp\gamma\sqrt{\log N/N}\to0$ & PROVED$_{\rm neg}$ & 123A\\
Startvektor-Resolvente $\to-1/z$ & PROVED$_{\rm neg}$ & H-T2\\
$b_1\to0\Rightarrow$ gesamter Limes diagonal & NO-GO & 123A\\
$b_{2,N}/b_{1,N}\to\infty$ & OPEN & 123F/G\\
abstraktes No-scalar-Lemma & CONDITIONAL & 123H\\
intrinsisches positives nichtskalares $W_N$ & OPEN & H-T3\\
gewichteter Rang-eins-Operator ist Projektor & NO-GO & 128A\\
NEU-128b Operator/Skalar-Formel & NO-GO & 128b\\
PSWF $\Rightarrow$ konkrete Prä-Lanczos-Metrik & HEURISTIC/OPEN & 130\\
Paper-VII-H3-Skalierung & NO-GO & 131\\
Cancellation $\Rightarrow$ absolute Schur-Zeilensumme & NO-GO & 131\\
algebraische $\beta$-Zerlegung & PROVED & 136\\
$\sum_{p\le N}(\log p)^2/p\sim\frac12(\log N)^2$ & PROVED & H-T4\\
Rohdivergenz aus Upper Bound & NO-GO & 136\\
$|c_p|^2=O((\log p)^2/p)$ intrinsisch & OPEN/CONDITIONAL & 134--135D\\
feste-$\beta$ $\Sigma_{\rm rel}^{\rm ren}\in\mathcal S_1$ & CONDITIONAL MODEL & 137\\
Fredholm-Basistheorie aus $\mathcal S_1$ & CONDITIONAL & 137/138\\
primeweise Eigenwert-/Eulerproduktlesart ohne T2 & NO-GO & 138\\
korrekte zweite Spur & PROVED & 139\\
intrinsisches T2 & OPEN & 142/143\\
$c_p\ne0$ für alle Primkanäle & OPEN & P05/H-T4\\
primdiagonales Mangoldt-$R$ & CONDITIONAL/OPEN & 140--144\\
gewöhnliche Mangoldt-Spur für $\Re\beta>1$ & CONDITIONAL MODEL & 141/144\\
$\Tr_{\rm reg}:=\AC[-\zeta'/\zeta]$ & DEFINITION & 145\\
Prime-only-Schichtzerlegung & PROVED & 146\\
Prime-only-Explizitformel ohne $\vartheta/\psi$-Korrektur & NO-GO & 147\\
RH-Verbindung des Prime-only-Defekts & PARTIAL/OPEN & 147\\
Mellin-Identität für $S_{\varphi,X}$ & NO-GO & 148\\
Mellin-Identität für $\Psi_{\varphi,X}$ & PROVED & H-T5\\
$\widehat\varphi$ ganz / $\widehat\varphi(0)=1$ & NO-GO & 148\\
$\operatorname{Res}_0\widehat\varphi=1$ & PROVED & 149\\
historische $\Psi-S$-Differenzformel & NO-GO & 148.6\\
quantitativer/uniformer $\Psi/S$-Transfer & OPEN & H-T5\\
fixed-contour Restlemma für $\Psi$ & CONDITIONAL & 149\\
uniforme Kontur und vollständige Residuen & OPEN & 149\\
Primlabel-Observable $N_{\mathbb P}$ & CONDITIONAL & 150\\
Primlabel-Spuridentität & CONDITIONAL MODEL & 150\\
Primlabel-Finite-Part $=-\zeta'/\zeta$ & OPEN & 150\\
operatorielle Realisierung von $\Tr_{\rm reg}$ & OPEN & 145/150\\
$R$-Cutoff-Transfer & OPEN & 146/150\\
\bottomrule
\end{longtable}
\normalsize

% ----------------------------------------------------------------
\section{Offene Kernfragen und Routing}
% ----------------------------------------------------------------

\begin{openproblem}[Root-Blocker nach P08]\label{op:root-blockers}
Offen bleiben insbesondere:
\begin{enumerate}
  \item konkrete KMS/GNS--Jacobi-Realisierung und Formkompatibilität;
  \item der strenge Nachweis $b_{2,N}/b_{1,N}\to\infty$;
  \item Konstruktion eines intrinsischen positiven nichtskalaren $W_N$;
  \item intrinsische Definition und quantitative Kontrolle von $|c_p|^2$;
  \item intrinsisches T2;
  \item Nichtentartung $c_p\ne0$;
  \item selbstadjungierte Mangoldt-$R$-Realisierung;
  \item uniforme Mellin-Kontur und vollständige Residuenzählung;
  \item der $\Psi/S$-Transfer in der für Finite Parts nötigen Stärke;
  \item Primlabel-Finite-Part und operatorielle Realisierung von $\Tr_{\rm reg}$;
  \item der quantitative $R$-Cutoff-Transfer;
  \item Grenzoperator-/Spektralmaßidentifikation des NEU-124-Pfads.
\end{enumerate}
\end{openproblem}

\begin{remark}[Routing]\label{rem:routing}
Gesicherte negative Resultate werden nach P10 geroutet. Intrinsische
Lift-/Quell-/Gramgeometrie, T2, Nichtentartung und globale
Fredholm-/Schattenfragen gehören nach P11. Finite-to-Infinite-Weil-Grenzfragen
werden nach P12 beziehungsweise spätere Grenzblöcke weitergereicht.
\end{remark}

% ----------------------------------------------------------------
\section{Provenienz und SYN-Endurteil}
% ----------------------------------------------------------------

\begin{remark}[Provenienz]\label{rem:provenance}
Dieses Manuskript konsolidiert den Live-Block
\texttt{04-grenzoperator-renormierung} gemäß
\texttt{audits/AUDIT-2026-08-09\_P08\_PassA\_FINAL\_SEAL.md}.
Die kanonische detaillierte Markdown-Quelle ist
\texttt{papers/P08\_Renormalized\_Prime\_Operators\_and\_Finite\_Part\_Structures.md}.
Die dortige Knotenprovenienz führt die 41 Live-Dokumente einschließlich der
beiden NEU-123F-Dateien. P05, P06 und P07 werden nur als bereits eingefrorene
Typ-/Spektral-/Herglotz-Firewalls benutzt.
\end{remark}

\begin{remark}[Schutzsatz]\label{rem:final-firewall}
Der belastbare P08-Kern ist
\[
  \boxed{
  \text{Jacobi-Kollapsdiagnose}
  \;\oplus\;
  \text{conditional feste-}\beta\text{-Spurklasse}
  \;\oplus\;
  \text{exakter Mangoldt-Mellin-Kanal}
  }
\]
mit offenen Operator- und Cutoff-Brücken. Insbesondere
\[
  \boxed{
  \text{analytische Fortsetzung von }-\zeta'/\zeta
  \ne
  \text{bereits konstruierter operatorieller Finite Part}
  \ne
  \text{Hilbert--Pólya-/Objekt-X-Operator}.
  }
\]
\end{remark}

\end{document}
